\chapter{Testování}
V této kapitole budou nejprve představeny nástroje, které byly využity při~vývoji aplikace a které připívají k větší kvalitě vytvářeného kódu. Dále bude přiblížen průběh uživatelského testování, jeho výsledky a zpětná vazba od~účastníků.

\section{Automatické testy a kontroly}
Všechny níže uvedené testy a kontroly jsou integrovány do CI/CD procesu. Spouští se automaticky před sloučením změn a při přidání commitu do hlavní vývojové větve. Díky tomu je minimalizováno riziko zanesení chyb do produkčního prostředí a je garantováno dodržování stanovených standardů.

\subsection{Larastan}
Pro kontrolu PHP kódu byl zvolen nástroj Larastan. Jedná se o rozšíření populárního statického analyzátoru PHPStan, které je optimalizováno pro framework Laravel. Larastan při procesu analýzy kódu spouští aplikační kontejner, což mu umožňuje kromě statické analýzy datových typů a logické struktury programu kontrolovat i datové typy, které jsou běžně zjistitelné až za běhu aplikace. I díky tomuto přístupu je zajištěna podpora pro mnoho dynamických vlastností frameworku Laravel, což výrazně přispívá k typové bezpečnosti a robustnosti aplikace. \cite{Larastan}

\subsection{Laravel Pint}
Konzistence formátování PHP kódu je zajišťována nástrojem Laravel Pint. Jde o nástroj postavený na základech PHP CS Fixer, který však nabízí zjednodušenou konfiguraci a sady předpřipravených pravidel odpovídajících standardům pro vývoj PHP aplikací. Pint byl v rámci projektu konfigurován tak, aby kontroloval a opravoval styl zápisu kódu dle standardu PSR-12. Zajištění jednotného stylu, zahrnujícího pravidla pro odsazování, závorkování či deklaraci typů, je klíčové pro čitelnost kódu a efektivní spolupráci v týmu. \cite{Pint}

\subsection{ESLint, Prettier}
V kontextu klientské části aplikace slouží pro statickou analýzu kódu nástroj ESLint. Jeho účelem je identifikace problematických míst v komponentách frameworku Vue.js a v kódu jazyka JavaScript, která by mohla vést k chybám za běhu aplikace nebo k nekonzistencím v logice komponent.

Jako doplněk k nástroji ESLint je pro formátování zdrojového kódu frontendu využíván Prettier. Zatímco ESLint se zaměřuje na kvalitu a logickou správnost kódu, Prettier zajišťuje, aby všechny skripty a Vue.js komponenty měly jednotné formátování bez ohledu na to, kým byly napsány.

\subsection{PHPUnit}
Pro automatizované testování funkčnosti aplikace je využíván testovací framework PHPUnit. V současné fázi vývoje obsahuje projekt sadu základních testů, které byly součástí využitého Vue starter kitu. Tyto testy ověřují klíčové procesy, jako je registrace nového uživatele, změna přístupového hesla či~dostupnost hlavního dashboardu po úspěšném přihlášení. V budoucnu by však bylo vhodné tyto testy rozšířit na další funkcionality systému jako je například synchronizace dat se systémy pro správu verzí a ověření správnosti zobrazovaných metrik na dashboardu.

\section{Uživatelské testování}
V této sekci bude popsán proces a výsledky uživatelského testování vytvořené aplikace. Hlavním cílem této fáze bylo ověřit, zda je navržené uživatelské rozhraní dostatečně intuitivní pro koncové uživatele, a identifikovat potenciální nedostatky, které by mohly negativně ovlivnit uživatelskou přívětivost. Na základě získaných poznatků byly následně navrženy úpravy systému.

Před zahájením samotného testování byli účastníci seznámeni se skutečností, že předmětem hodnocení je samotný systém, nikoliv jejich individuální schopnosti či dovednosti. Testování ze zúčastnili celkem 4 uživatelé a bylo realizováno prezenční formou. To umožnilo sledovat bezprostřední reakce uživatelů a prováděné akce při plnění úkolů.

\subsection{Průběh testování}
Samotný proces testování byl rozdělen do několika fází:

\begin{enumerate}
    \item \textbf{Seznámení se systémem:} Každému účastníkovi byla nejprve představena podstata aplikace a její hlavní účel.
    \item \textbf{Průchod testovacími scénáři:} Uživatelům byly předány připravené scénáře obsahující sadu konkrétních úkolů, které měli v systému samostatně vykonat.
    \item \textbf{Závěrečný rozhovor:} Po dokončení všech úkolů byl s každým uživatelem veden strukturovaný rozhovor. Během této diskuse byly shromážděny subjektivní poznatky a náměty na funkční či vizuální vylepšení.
\end{enumerate}

\subsection{Testovací scénáře}

\subsubsection{Registrace do systému}
\begin{description}
    \item[\textcolor{black}{Modelová situace}] \hfill \\
        Dozvěděli jste se o nástroji, který umí porovnávat aktivitu vývojářů. Jste zvědaví, jak si stojíte oproti svým kolegům.
    \item[\textcolor{black}{Úkol}] \hfill
        \vspace{-2mm}
        \begin{enumerate}
            \item Zaregistrujte se do aplikace pod jménem Tomáš Novák, s e-mailem \mbox{tnovak@gmail.com} a libovolným heslem.
        \end{enumerate}
\end{description}

\subsubsection{Změna nastavení}
\begin{description}
    \item[\textcolor{black}{Modelová situace}] \hfill \\
        Při registraci jste omylem udělal ve své e-mailové adrese překlep a chtěl byste to napravit. Zároveň byste si chtěl změnit barevné schéma aplikace.
    \item[\textcolor{black}{Úkoly}] \hfill
        \vspace{-2mm}
        \begin{enumerate}
            \item Upravte svou e-mailovou adresu na t.novak@gmail.com
            \item Změňte vzhled aplikace na jiné barevné schéma.
        \end{enumerate}
\end{description}

\subsubsection{Zobrazení metrik}
\begin{description}
    \item[\textcolor{black}{Modelová situace}] \hfill \\
        Jste vedoucí vývojového týmu a chcete si zobrazit informace o aktivitách svého kolegy.
    \item[\textcolor{black}{Úkol}] \hfill
        \vspace{-2mm}
        \begin{enumerate}
            \item Zobrazte si metriky uživatele Lukáš Paukert od 20. ledna do 29. ledna.
        \end{enumerate}
\end{description}

\subsubsection{Správa sledovaných repozitářů}
\begin{description}
    \item[\textcolor{black}{Modelová situace}] \hfill \\
        Jste vedoucí vývojového týmu a chcete upravit seznam repozitářů, u kterých dochází ke sběru dat.
    \item[\textcolor{black}{Úkoly}] \hfill
        \vspace{-2mm}
        \begin{enumerate}
            \item Začněte od dnešního dne evidovat aktivity z repozitáře laravel/laravel z veřejné instance GitHub. Repozitář by se měl v systému zobrazovat pod jménem Laravel a data by měla být synchronizována každých 6~hodin.
            \item Odeberte repozitář Yii2 ze sledovaných projektů (včetně všech jeho dat).
        \end{enumerate}
\end{description}

\subsubsection{Přiřazení účtu ze systému pro správu verzí}
\begin{description}
    \item[\textcolor{black}{Modelová situace}] \hfill \\
        Jste vedoucí vývojového týmu a chcete novému uživateli přiřadit účet z verzovacího systému.
    \item[\textcolor{black}{Úkol}] \hfill
        \vspace{-2mm}
        \begin{enumerate}
            \item Přiřaďte uživateli Lukáš Paukert účet se jménem lukas.paukert z verzovacího systému GitLab provozovaným firmou DAMI.
        \end{enumerate}
\end{description}

\subsection{Průběh a vyhodnocení}
